%% kapa: exam=1 round=19 year=2008 period=1 form=A subject=civil_law subject_ko=민법 seq=39 number=39
%% source: Q-Net 감정평가사 기출문제 (cst00309) artlSeq=1001962 file=extracted/1008/1001962/2008년 제19회 감정평가사 1차시험 기출문제/1교시.hwp
%% conversion: kordoc 4.24.51 -> tools/gen1.py (auto); review=proofread-sonnet; flags=none
\begin{question}{39}
저당권에 관한 설명으로 옳은 것을 모두 고른 것은? (다툼이 있는 경우에는 판례에 의함)
\begin{framedtext}
ㄱ. 불법말소된 저당권등기가 회복되기 전에 경매가 행하여져 매수인이 매각대금을 완납하였다면 저당권말소등기의 회복등기를 청구할 수 없다. \par
ㄴ. 채무자의 변제로 피담보채권이 소멸하면 말소등기를 하지 않아도 저당권은 소멸한다. \par
ㄷ. 저당권의 효력은 저당부동산에 부합된 물건에 미치므로, 명인방법을 갖춘 수목에도 토지저당권의 효력이 미친다. \par
ㄹ. 토지에 관하여 저당권이 설정될 당시 그 지상에 존재하는 건물이 미등기상태였다면 그 건물을 위한 법정지상권이 성립할 수 없다. \par
ㅁ. 근저당권이전의 부기등기가 경료된 경우, 피담보채무의 소멸을 원인으로 한 근저당권설정등기 말소청구의 상대방은 양수인이 아니라 양도인이다.
\end{framedtext}
\begin{choices}
  \item ㄱ, ㄴ
  \item ㄴ, ㅁ
  \item ㄱ, ㄷ, ㅁ
  \item ㄱ, ㄹ, ㅁ
  \item ㄴ, ㄷ, ㄹ
\end{choices}
\end{question}
\answer{1}
